\section{Reproducibility}
\label{sec:repro}

\subsection{Software and hardware}
Experiments ran on a 24-core workstation shared with unrelated jobs, under
PennyLane 0.42.3, PyTorch 2.6.0, NumPy 1.26.4 and pandas 2.2.3 on Python
3.12.9. Quantum circuits were simulated on the \texttt{default.qubit}
state-vector backend with backpropagation. All reported models are trained on
CPU: for an eight-qubit register the simulator is CPU-bound, and a GPU gives no
benefit.

\paragraph{Measurement under contention.}
Wall-clock figures are sensitive to what else the machine is running. The same
training step measured \SI{16.4}{\second} at a one-minute load average of 57 on
24 cores and \SI{0.38}{\second} once the machine was scheduled, a factor of 43.
Every timing reported here was taken through a job scheduler that admits work
against declared CPU, RAM and VRAM, and the load at measurement time is
recorded alongside the figure.

\subsection{Determinism}
Each (model, fold, seed) cell seeds PyTorch, NumPy and the training-set
subsample from the same integer. Circuit parameters are initialised from a
dedicated generator so that the \textsc{Frozen} condition draws exactly the same
weights as the \textsc{Trained} condition's initialisation. Results are
checkpointed per cell, so an interrupted sweep resumes without recomputing
completed work and without changing any seed.

\subsection{Data and code}
The dataset is the public JPX Tokyo Stock Exchange
panel~\citep{jpx2022competition}. The analysis code, including the evaluation
statistics, the walk-forward protocol and the batched circuit, is released with
the paper. The calibration checks reported in
Section~\ref{sec:protocol:stats} --- in particular that the
backtest-overfitting estimator returns \num{0.53(4)} on independent random
strategies --- are part of the test suite.

\subsection{Deviations from the reference implementation}
The pipeline this work audits --- our own, as stated in
Section~\ref{sec:intro} --- differs in four respects that materially affect its
numbers, each corrected here: the split adjustment direction
(Section~\ref{sec:data:adjustment}); the Sharpe ratio, which the released code
computes as mean over standard deviation with no annualisation, though the
published paper states that its metric is annualised; the evaluation window,
which date subsampling reduced to \num{22} trading days; and the absence of
transaction costs. Section~\ref{sec:results:reeval} reports the original models
re-scored over \num{241} days with all four corrected, so that the difference
is attributable rather than asserted.
